\section{Conclusion}
\label{sec:conclusion}

Code-based phenotyping has been refined methodologically through
latent-class fits, anchor learning, semi-supervised algorithms, and
high-throughput pipelines, yet the structural reason its prevalence
estimates are biased has stayed informal. The masked-data framing
gives the structure a name. Informative coding is a C2 violation;
chart review is a singleton candidate set; the Rogan-Gladen plug-in is
the deconvolution it enables; marginal-fit blindness is the
identifiability cost of working with only the cohort code frequency;
the case-mix gap is the limit on what review can recover.

This recasting moves the practitioner's question from method-specific
(``should I use PheCAP or PheNorm?'') to assumption-specific (``does
my coding mechanism violate C2, and is my chart-reviewed subsample
representative of the cohort?''). The diagnosis is almost always that
C2 fails. The framework then says precisely what is required: singleton
candidate sets, supplied by chart review, anchor observations, or a
gold-standard registry, sampled to match the cohort case mix; a
reported case-mix gap as a sensitivity quantity; and recognition that
the marginal-fit check is silent on prevalence error.

The same coarsening structure appears across health-data measurement
problems with the same informative-observation issue: claims-based
comorbidity indices, registry case-finding, cause-of-death coding, and
adverse-event surveillance. In each case a latent clinical truth is
observed only through an administrative process whose intensity depends
on the truth, and the C1-C2-C3 conditions and the identifiability
machinery transfer directly.

Three extensions sit naturally next. First, multi-code and temporal
candidate sets: real phenotype algorithms aggregate code counts, code
timing, medication sequences, and lab trajectories, and the
identifiability rank conditions extend but have not been worked out.
Second, chart-review error: adjudication is itself imperfect, and the
C2-relaxation results of \citet{towell2026mdrelax} would let the
singleton assumption itself be relaxed. Third, propagation: code-based
phenotypes feed downstream association and causal-inference analyses,
where informative coding induces differential misclassification of the
exposure or outcome, and the propagation of the coarsening bias into
those estimands is open.

The contribution this paper makes is positioning. It places electronic
phenotyping inside a unified vocabulary that names the failing
assumption (C2), supplies the remedy (singleton candidate sets), and
quantifies the residual (the case-mix gap). The technical results are
not breakthroughs in isolation, they are corollaries of the
masked-data framework applied to the specific structure of EHR coding,
but the unification is what lets a clinician or methodologist read off,
from the coding mechanism and the chart-review sampling, the direction
and bounded magnitude of the prevalence error they will incur.
